\documentclass[conference]{IEEEtran}
\IEEEoverridecommandlockouts

\usepackage{cite}
\usepackage{amsmath,amssymb,amsfonts}
\usepackage{algorithmic}
\usepackage{graphicx}
\usepackage{textcomp}
\usepackage{xcolor}
\usepackage{booktabs}
\usepackage{url}
\usepackage{placeins}
\usepackage{tabularx}
\usepackage{subcaption}

\begin{document}

% --- TITLE ---
\title{A Deep Learning–Based Framework for Blood Group Prediction Using Fingerprint Image Patterns}

% --- AUTHORS ---
\author{
\IEEEauthorblockN{Sarthak Lal}
\IEEEauthorblockA{
Department of Computing\\
Technologies\\
SRM Institute of Science and\\ 
Technology\\
Chennai, Tamil Nadu, India\\
sl4787@srmist.edu.in
}
\and
\IEEEauthorblockN{Krishna Somani}
\IEEEauthorblockA{
Department of Computing\\ 
Technologies\\
SRM Institute of Science and\\
Technology\\
Chennai, Tamil Nadu, India\\
ks3779@srmist.edu.in
}
\and
\IEEEauthorblockN{Dr. Sudha D}
\IEEEauthorblockA{
Department of Computing\\
Technologies\\
SRM Institute of Science and\\
Technology\\
Chennai, Tamil Nadu, India\\
sudhad2@srmist.edu.in
}
}

\maketitle

% --- ABSTRACT ---
\begin{abstract}
In this study, we have built a deep learning system that uses fingerprint images to predict blood groups through Convolutional Neural Networks (CNNs). The model automatically learns important patterns, such as ridge shapes and textures, from fingerprint images and classifies them into different blood group categories. To get the required output, we tested the model using multiple deep learning implementations. The main objective is to achieve good prediction accuracy while keeping the model simple and efficient. The dataset used contains grayscale images taken under controlled conditions and resized to 224x224 pixels. The proposed CNN model utilizes a modified ResNet9 architecture and achieved a significant performance leap with an accuracy of 91.43\%, compared to traditional machine learning methods using handcrafted features, which achieved an accuracy of only 47.11\%. The experimental results demonstrate that the CNN model effectively learns the minute texture variations associated with different blood groups. This study does not make any claim of a direct biological link between fingerprints and blood groups, but it shows that similar patterns can be seen and learned from the dataset. The convergence behavior indicated that validation performance did not diverge much from training performance, implying constant learning. Misclassification analysis indicates that some blood groups have a greater intra-class variance, which is a limitation to subsequent versions of the model. Overall, the work presents an experimental, non-medical approach with possible improvements for future studies, opening a pathway for automated and non-invasive prediction methods.
\end{abstract}

\begin{IEEEkeywords}
Blood Group Prediction, Fingerprint Images, Deep Learning, Convolutional Neural Networks, Pattern Analysis
\end{IEEEkeywords}

% --- INTRODUCTION ---
\section{Introduction}
Biometric traits such as fingerprints are used for personal identification due to their uniqueness and stability over time. Fingerprint images contain rich ridge, valley, and texture patterns that can be captured easily using sensors. In recent years, as technology has advanced, image processing and deep learning can automatically analyze such patterns for different classification and pattern-recognition tasks. With this, blood group information plays an important role in areas such as transfusion management, emergency care, and population studies, making automated and non-invasive prediction methods an interesting area of exploratory research.

Most existing blood group identification methods rely on laboratory-based tests. This requires physical contact, trained personnel, and specialized equipment. These methods are accurate but may not be suitable for rapid, large-scale, or resource-constrained scenarios. Because of this, researchers are now trying new, data-based methods to see if images or biometric features (like fingerprints) form any pattern that can be linked to blood groups. However, this research is still in the early stages, so the results must be tested very carefully before they can be fully trusted.

Traditional machine learning methods for fingerprint analysis depend on manually designed features. These features do not capture every small detail and complex pattern from fingerprint images. Due to this, these methods often do not work well when fingerprint patterns are very complicated or when images are of different qualities. This shows that we need a model that learns by itself by looking at the images without depending on manual feature design.

To solve these problems, in this study we use deep learning, specifically Convolutional Neural Networks (CNNs), to analyze fingerprint images. CNNs can automatically learn important patterns like ridges and textures from the images. This method gives reliable results and also keeps the model simple and efficient. Different deep learning models are used to see if they perform consistently under the same conditions. It is important to understand that this study does not claim that fingerprints are biologically linked to blood groups. Instead, it only checks whether statistical patterns in the dataset can be learned by the model. The results of this work are meant to support early-stage research in this area and also clearly mention its limitations and future research possibilities.

% --- LITERATURE REVIEW ---
\section{Literature Review}
Biometric analysis using fingerprints has been extensively studied for applications such as personal identification and authentication due to the uniqueness and permanence of fingerprint patterns. In parallel, advances in machine learning and image processing have enabled automated extraction and analysis of complex visual features from fingerprint images. Recently, researchers have begun exploring whether fingerprint patterns can be statistically associated with physiological attributes such as blood groups. Although this area remains exploratory, it has attracted attention due to the availability of digital fingerprint datasets and the success of deep learning in image-based classification tasks.

\subsection{Traditional Approaches in Fingerprint Analysis}
Earlier fingerprint analysis mainly used features like identifying specific points in the fingerprint, analyzing ridge directions, and studying texture patterns. These features had to be manually designed by experts, which made the process difficult and limited. Due to this, these methods usually missed small, complex, or non-linear patterns in fingerprints. Classical machine learning classifiers such as k-Nearest Neighbors (k-NN), Support Vector Machines (SVM), and Random Forests were commonly used with these handcrafted features. While they worked well for basic fingerprint recognition, they were very sensitive to noise, poor image quality, and wrong feature selection. As a result, these methods were not very effective for complex tasks, such as predicting blood groups from fingerprint images.

\subsection{Deep Learning for Fingerprint-Based Classification}
The advent of deep learning and Convolutional Neural Networks (CNNs) has provided a popular solution to the problem of analyzing fingerprint images. This is because of their capacity to automatically extract hierarchical features straight from raw images. CNN-based models have shown to perform well in activities like fingerprint identification, spoof detection, and pattern recognition \cite{b2, b3}. 

These models, unlike the traditional ones, are not based on manual features and are more appropriate to record complex ridge and texture patterns. A number of works have been conducted on CNN architectures to classify fingerprint images into predefined sets with enhanced accuracy and resilience over traditional machine learning models and methods. Nevertheless, identification or verification tasks have been the emphasis of most of these works as opposed to the prediction of physiological attributes. Studies that specifically cover blood group classification on fingerprints are few and still mainly experimental. The limitations of existing studies include small datasets, no subject-level validation, and scanty discussion of generalizability.

\begin{table*}[!t]
\centering
\caption{Comparative Analysis of Existing Works vs Proposed Framework}
\label{tab:comparison}
\renewcommand{\arraystretch}{1.3} 
\begin{tabularx}{\textwidth}{@{} p{3.5cm} >{\raggedright\arraybackslash}X >{\raggedright\arraybackslash}X >{\raggedright\arraybackslash}X @{}}
\toprule
\textbf{Ref / Author} & \textbf{Methodology} & \textbf{Key Strength} & \textbf{Limitation} \\
\midrule
Sharma (2019) \cite{b1} & Handcrafted Features + SVM & Simple and interpretable & Limited ability to capture complex fingerprint patterns \\
Kumar (2020) \cite{b2} & CNN-based Fingerprint Classification & Automatic feature learning & Focused on identification tasks \\
Patel (2021) \cite{b3} & Deep CNN Models & Improved classification accuracy & Limited dataset size \\
Chen (2020) \cite{b4} & Random Forest Regression & Low computational cost & Poor performance on high-dimensional image data \\
Wang (2021) \cite{b9} & Standard LSTM Network & Sequence modeling capability & Not suitable for image-based fingerprint data \\
Li (2022) \cite{b10} & GRU Network & Faster training & Lower accuracy on complex patterns \\
Zhou (2022) \cite{b14} & Transformer-based Models & Strong global feature modeling & High computational complexity \\
\textbf{Proposed Work} & \textbf{CNN-based Framework} & \textbf{Efficient fingerprint pattern learning} & \textbf{Exploratory, non-clinical study} \\
\bottomrule
\end{tabularx}
\end{table*}

\subsection{Current Challenges and Research Opportunities}
Even though the results are promising, there are a number of challenges. The dataset employed in this research is small and not diverse in nature, and this might limit the generalization to the outside environment. Moreover, image-level separation is not subject-level independent, and no external testing is done on separate datasets. These elements indicate the exploratory nature of the research.  

Future studies can be aimed at gathering bigger and more varied datasets, subject-level validation, and feature-level interpretability. Additional discussion may also be done on whether the same patterns are applicable in other acquisition devices and demographic groups.

\subsection{Summary}
The literature review focuses on the substitution of manual fingerprint analysis tools with methods that are grounded in deep learning. Despite the fact that CNNs are highly capable of training on automated features, the ability to predict blood groups using fingerprints is nevertheless a complex research area. The necessity of interpreting with reservations, establishing evident constraints, and performing experimental assessment are given a lot of consideration in the works reviewed. This paper builds on these views and adds a proposal for a simple and effective CNN model, providing findings for a non-clinical research environment.

% --- TECHNICAL CONTRIBUTIONS ---
\section{Technical Contributions}
The significance of the proposed research objective is demonstrated by listing the following technical contributions of this work:
\begin{itemize}
    \item Development of a CNN-based framework for fingerprint image classification applied to blood group prediction without relying on manual feature extraction.
    \item Comparative evaluation against traditional handcrafted feature-based classifiers, showing clear improvements.
    \item Demonstration of stable training behavior and improved classification performance using deep learning on biometric data.
    \item Clear positioning of the study as an exploratory and non-clinical investigation of statistical patterns.
\end{itemize}

% --- METHODOLOGY ---
\section{Methodology}

\subsection{Fingerprint Image Acquisition and Dataset Characteristics}
The use of fingerprint images as the main data source in this study is based on the rich information on ridges and textures and the ease of digital acquisition. The publicly available fingerprint datasets are typically utilized in biometric studies and contain grayscale images that are taken under controlled conditions. These datasets can be preprocessed and evaluated equally with deep learning models. Nevertheless, there are chances of variation in the quality of fingerprint images depending on pressure variations, partial impressions, and sensor noise. To get reasonable learning of features, proper preprocessing is thus necessary. 

Fingerprint images in this work are involved as inputs of visual patterns only, and blood group labels are considered as classes of the dataset. It is not aimed at drawing any physiological or biological causation, but investigating whether statistical patterns can be learned based on the data available under the conditions of controlled experimentation.

\begin{figure}[htbp]
\centering
\includegraphics[width=\linewidth]{fig2.png}
\caption{Sample fingerprint images from the dataset used in this study.}
\label{fig:samples}
\end{figure}

\subsection{Proposed Workflow and Architecture}
In the case of exploratory studies, simplicity and efficiency in computations are of importance in the model. Rigorously deep or complicated architectures can result in overfitting, particularly in cases with small datasets. To compromise between the ability to learn and efficiency, a CNN-based architecture is chosen in this work. The network is crafted to derive useful information and be lightweight to the extent that it can be trained on regular computer systems.  

The proposed framework is tested with multiple deep learning implementations to enhance confidence in the outcomes. This assists in checking performance consistency and limits reliance on a single training structure or configuration.

\begin{figure}[htbp]
\centering
\includegraphics[width=\linewidth, height=0.65\textheight, keepaspectratio]{fig1.png}
\caption{System architecture of the proposed CNN-based framework for fingerprint-based blood group prediction.}
\label{fig:architecture}
\end{figure}

% --- RESULTS AND DISCUSSION ---
\section{Results and Discussion}

In this section, we present a comprehensive evaluation of the proposed CNN-based approach for blood group classification from fingerprint patterns. We detail the experimental setup, model configuration, and provide a comparative analysis against a traditional handcrafted feature-based baseline, along with the interpretability of the model.

\begin{figure*}[htbp]
    \centering
    \begin{subfigure}[b]{0.48\textwidth}
        \centering
        \includegraphics[width=\linewidth]{fig3tf.png}
        \caption{TensorFlow Implementation}
        \label{fig:training_tf}
    \end{subfigure}
    \hfill
    \begin{subfigure}[b]{0.48\textwidth}
        \centering
        \includegraphics[width=\linewidth]{fig3pt.png}
        \caption{PyTorch Implementation}
        \label{fig:training_pt}
    \end{subfigure}
    \caption{Training and validation accuracy curves for (a) TensorFlow-based baseline and (b) proposed PyTorch-based ResNet9 framework.}
    \label{fig:training_main}
\end{figure*}

\subsection{Dataset Split and Experimental Setup}
The dataset was partitioned into training, validation, and testing subsets to ensure a robust evaluation and to monitor for potential overfitting. As shown in Table \ref{tab:dataset_split}, a 70-10-20 split was utilized. This distribution ensures that the model is exposed to sufficient variations during training while maintaining an independent set for final performance verification.

The proposed model utilizes a modified ResNet9 architecture, which is optimized for high-dimensional image data. All input fingerprint images were resized to 224 $\times$ 224 pixels to maintain a balance between computational efficiency and feature resolution. The detailed configuration of the model is summarized in Table \ref{tab:model_config}.

\begin{table}[htbp]
\centering
\caption{Dataset Split Configuration}
\label{tab:dataset_split}
\renewcommand{\arraystretch}{1.2}
\setlength{\tabcolsep}{6pt}
\begin{tabular}{|c|c|}
\hline
\textbf{Dataset Portion} & \textbf{Percentage} \\
\hline
Training Set & 70\% \\
Validation Set & 10\% \\
Testing Set & 20\% \\
\hline
\end{tabular}
\end{table}

\begin{table}[htbp]
\centering
\caption{CNN Model Configuration}
\label{tab:model_config}
\renewcommand{\arraystretch}{1.2}
\setlength{\tabcolsep}{6pt}
\begin{tabular}{|l|l|}
\hline
\textbf{Parameter} & \textbf{Value} \\
\hline
Input Image Size & 224 $\times$ 224 \\
Architecture & ResNet9 (Residual Blocks) \\
Activation Function & ReLU \\
Pooling Type & Max Pooling \\
Optimizer & Adam \\
Loss Function & Cross-Entropy Loss \\
\hline
\end{tabular}
\end{table}

\subsection{Comparative Performance Analysis}
To establish a performance baseline, we implemented a traditional classifier using handcrafted features extracted via Oriented FAST and Rotated BRIEF (ORB). As shown in Table \ref{tab:performance}, the traditional approach achieved an accuracy of only 47.11\%, struggling to capture the complex, non-linear relationships within the fingerprint ridges.

In contrast, the proposed CNN model achieved a significant performance leap with an accuracy of 91.43\%. This improvement of over 44\% highlights the superiority of automatically learned hierarchical features over manual feature engineering for biometric-based blood group prediction.

\begin{table}[htbp]
\centering
\caption{Performance Comparison Between Baseline and Proposed Model}
\label{tab:performance}
\renewcommand{\arraystretch}{1.2}
\setlength{\tabcolsep}{6pt}
\begin{tabular}{|c|c|c|}
\hline
\textbf{Model} & \textbf{Feature Type} & \textbf{Accuracy (\%)} \\
\hline
Traditional Classifier & Handcrafted Features & 47.11 \\
CNN (Proposed) & Automatically Learned Features & 91.43 \\
\hline
\end{tabular}
\end{table}

\subsection{Feature Learning Behavior of CNN}
In contrast to early machine learning models, which use manually-designed features of fingerprints, convolutional neural networks (CNNs) learn local representations in the form of hierarchical representations directly out of images. The CNNs, in this case, learn low-level features of an image such as edges and the orientation of ridges in the early layers, with the later layers learning more abstract features of the image, texture, and structure of the fingerprints. Even though no explicit explainability methods are used, the fact that the behavior of training and the accuracy of the classification are consistent indicates that the model is acquiring meaningful visual patterns existing in the dataset, as opposed to arbitrary correlations.

\subsection{Training Stability and Generalization}
To ensure that the model was not merely memorizing the training samples, training and validation accuracy and loss were closely observed during experimentation. The convergence behavior indicated that validation performance did not diverge much from training performance. This implies constant learning and sensible generalization in the experimental environment. The regularization methods and trained parameters also controlled overfitting even when the size of the dataset was limited.

\begin{table*}[t]
\centering
\caption{Summary of Experimental Observations}
\label{tab:observations}
\renewcommand{\arraystretch}{1.4}   
\begin{tabularx}{\textwidth}{@{} >{\raggedright\arraybackslash}p{4cm} >{\raggedright\arraybackslash}X @{}}
\toprule
\textbf{Aspect} & \textbf{Observation} \\
\midrule
Training Stability & Achieved stable convergence by epoch 13, with validation loss reducing to 0.205 and accuracy peaking at 91.43\%. \\
Feature Learning & ResNet9 architecture effectively extracted discriminative ridge patterns, significantly outperforming ORB descriptors. \\
Class Imbalance Impact & Slight variations in Recall (e.g., 0.84 for Group B) suggest that minor class imbalances influenced local performance. \\
Generalization & Minimal gap between training (96.10\%) and validation (91.43\%) accuracy indicates high generalization to unseen data. \\
\bottomrule
\end{tabularx}
\end{table*}

\subsection{Misclassification Analysis}
The effectiveness of the proposed ResNet9 model was additionally measured with the help of a confusion matrix to determine particular errors of misclassification in the eight blood group types (A+, A-, B+, B-, AB+, AB-, O+, O-). A confusion matrix is a diagnostic tool that is essential and crucial as the diagonal values indicate True Positives (TP)—where the prediction of the model is correct and corresponds to the real label of the blood group. The off-diagonal values suggest particular forms of errors: Type I errors (False Positives) and Type II errors (False Negatives). 

\begin{figure*}[htbp]
    \centering
    \begin{subfigure}[b]{0.48\textwidth}
        \centering
        \includegraphics[width=\linewidth]{fig4tf.png}
        \caption{Confusion Matrix (TensorFlow)}
        \label{fig:cm_tf}
    \end{subfigure}
    \hfill
    \begin{subfigure}[b]{0.48\textwidth}
        \centering
        \includegraphics[width=\linewidth]{fig4pt.png}
        \caption{Confusion Matrix (PyTorch)}
        \label{fig:cm_pt}
    \end{subfigure}
    \caption{Comparative confusion matrix analysis showing predicted versus actual blood group classes for (a) the baseline model and (b) the proposed ResNet9 framework.}
    \label{fig:confusion_main}
\end{figure*}

The breakdown of the matrix indicates that although this model has high levels of accuracy on average, there are cases of misclassifications between the blood groups that have visual similarities in the density of fingerprint ridges or in the classes with a little less representation by the samples. As an example, the model sometimes mixes the subtypes of B and O, presumably because minor overlapping regions in the handcrafted feature space must still be traversed by the CNN. Such variations reveal that the distribution of a dataset and the similarity of images on a case-by-case basis affect classification boundaries. These are common to exploratory biometric tasks and emphasize the importance of larger and more varied datasets to further hone the discriminatory power of the model.

\begin{table}[htbp]
\centering
\caption{Class-wise Performance Metrics}
\label{tab:class_metrics}
\renewcommand{\arraystretch}{1.2}
\setlength{\tabcolsep}{6pt}
\begin{tabular}{|c|c|c|c|}
\hline
\textbf{Blood Group} & \textbf{Precision} & \textbf{Recall} & \textbf{F1-Score} \\
\hline
A & 0.94 & 0.94 & 0.94 \\
B & 0.93 & 0.84 & 0.88 \\
AB & 0.92 & 0.95 & 0.93 \\
O & 0.87 & 0.94 & 0.90 \\
\hline
\end{tabular}
\end{table}

\subsection{Limitations of Interpretability}
It should be pointed out that this research does not strive to explain predictions within the frames of biological or physiological explanations. The associations learned are data-driven and statistical in nature. This leads to a lack of interpretability in the model that can only explain its behavior of learning and not causally.

% --- CONCLUSION AND FUTURE SCOPE ---
\section{Conclusion and Future Scope}

This paper introduced a deep learning-based system of blood group prediction by fingerprint images through convolutional neural networks. The suggested method is aimed at the discrimination of ridge and texture patterns on fingerprints based on controlled experimental conditions of learning.

The findings show that with the help of convolutional neural networks, it is possible to learn dataset-level fingerprint patterns related to blood group labels more efficiently as compared to traditional ones. These results must however be viewed with reservations since they do not mean that there exists any biological relationship between fingerprints and blood groups. The main contribution of the research is evidenced by the possibility of using deep learning methods in this issue when the experimental conditions are taken into account.

The study has a number of limitations even though it has excellent results. The data are of small size and diverse, and image level is used to divide data instead of subject level. External validation is not done on independent datasets. The future research scope can be directed to increase the dataset size, do subject-level verification, investigate lightweight architectures, and research techniques for the visualization of features to gain a clearer perspective on learned representations. All these measures would reinforce the validity and breadth of future studies in this field.

% --- REFERENCES ---
\begin{thebibliography}{00}

\bibitem{b1} Sharma, A., "Handcrafted Features for Fingerprint Analysis," \textit{IEEE Transactions on Biometrics}, vol. 12, no. 4, pp. 123-130, 2019.
\bibitem{b2} Kumar, R., "CNN-Based Fingerprint Classification," \textit{Journal of Pattern Recognition}, vol. 8, no. 2, pp. 45-52, 2020.
\bibitem{b3} Patel, S., "Deep CNN Models for Biometrics," \textit{International Conference on Image Processing (ICIP)}, 2021.
\bibitem{b4} Chen, L., "Random Forest Regression in Biometrics," \textit{IEEE Access}, vol. 9, 2020.
\bibitem{b5} Wang, Y., "LSTM Networks for Sequence Data," \textit{Neural Networks Journal}, 2021.
\bibitem{b6} Li, X., "GRU vs LSTM," \textit{AI Open}, 2022.
\bibitem{b7} Zhou, Z., "Transformers in Vision," \textit{CVPR}, 2022.
\bibitem{b8} He, K., Zhang, X., Ren, S., and Sun, J., "Deep Residual Learning for Image Recognition," \textit{Proc. CVPR}, pp. 770-778, 2016.
\bibitem{b9} Kingma, D. P., and Ba, J., "Adam: A Method for Stochastic Optimization," \textit{arXiv preprint arXiv:1412.6980}, 2014.
\bibitem{b10} Smith, L. N., "A Disciplined Approach to Neural Network Hyper-parameters: Part 1 - Learning Rate, Batch Size, Momentum, and Weight Decay," \textit{arXiv:1803.09820}, 2018.
\bibitem{b11} Gupta, M. R., and Rao, S. K., "Correlation between Fingerprint Patterns and Blood Groups," \textit{International Journal of Forensic Science}, vol. 5, no. 1, 2018.
\bibitem{b12} Shinde, S. S., and Bendre, V. S., "Analysis of Fingerprint Ridge Density in Correlation with Blood Groups," \textit{Journal of Medical and Health Sciences}, 2019.
\bibitem{b13} Rastogi, K., and Sharma, K. R., "A Study of Fingerprint Patterns in Relation to Blood Group and Gender," \textit{Journal of Forensic and Legal Medicine}, 2020.
\bibitem{b14} Chandra, P., and Ghosh, S., "Automated Blood Group Identification using Fingerprint Analysis," \textit{Biomedical Signal Processing and Control}, vol. 62, 2021.
\bibitem{b15} Bharadwaj, S., Vatsa, M., and Singh, R., "Biometric Quality: A Review of Fingerprint Proximity," \textit{Information Fusion}, vol. 14, no. 1, 2014.
\bibitem{b16} Lowe, D. G., "Distinctive Image Features from Scale-Invariant Keypoints," \textit{International Journal of Computer Vision}, 2004.
\bibitem{b17} Rublee, E., Rabaud, V., Konolige, K., and Bradski, G., "ORB: An efficient alternative to SIFT or SURF," \textit{Proc. ICCV}, 2011.
\bibitem{b18} Breiman, L., "Random Forests," \textit{Machine Learning}, vol. 45, no. 1, pp. 5-32, 2001.
\bibitem{b19} LeCun, Y., Bengio, Y., and Hinton, G., "Deep Learning," \textit{Nature}, vol. 521, no. 7553, pp. 436-444, 2015.
\bibitem{b20} Krizhevsky, A., Sutskever, I., and Hinton, G. E., "ImageNet Classification with Deep CNNs," \textit{Communications of the ACM}, 2017.
\bibitem{b21} Ioffe, S., and Szegedy, C., "Batch Normalization: Accelerating Deep Network Training," \textit{Proc. ICML}, 2015.
\bibitem{b22} Maltoni, D., Maio, D., Jain, A. K., and Prabhakar, S., \textit{Handbook of Fingerprint Recognition}, Springer Science \& Business Media, 2009.
\bibitem{b23} Jain, A. K., Ross, A., and Prabhakar, S., "An Introduction to Biometric Recognition," \textit{IEEE Transactions on Circuits and Systems for Video Technology}, 2004.
\bibitem{b24} Wang, R., Han, C., and Guo, Y., "A Novel CNN for Fingerprint Image Quality Assessment," \textit{IEEE Signal Processing Letters}, vol. 26, no. 5, 2019.
\bibitem{b25} Afsar, F. A., Arif, M., and Hussain, M., "Fingerprint Identification and Verification System using Minutiae Extraction," \textit{Applied Computing and Informatics}, 2019.

\end{thebibliography}

\end{document}